% Vietnamese translation for OI Public Library
% Source-PDF: IOI中国国家候选队论文/2009/骆可强/论程序底层优化的一些方法与技巧.pdf
\documentclass[11pt]{article}
\usepackage{oipl-vn}

\newcommand{\Oh}{\mathcal{O}}
\newcommand{\RISC}{\textsc{Risc}}
\newcommand{\CISC}{\textsc{Cisc}}
\newcommand{\SIMD}{\textsc{Simd}}

\title{Bàn về một số phương pháp và kỹ thuật tối ưu tầng thấp của chương trình}
\author{Luo Keqiang\\Trường Trung học số 7 Thành Đô}
\date{Luận văn Trại đông 2009}

\begin{document}
\maketitle

\origtitle{Bàn về một số phương pháp và kỹ thuật tối ưu tầng thấp của chương trình}
\sourcepath{IOI China candidate team papers / 2009 / Luo Keqiang}

\begin{abstract}
Bài viết lấy mục tiêu tối ưu hiệu quả thời gian của chương trình, nhìn từ các
khái niệm tương đối tầng thấp như trình biên dịch, mã hợp ngữ và đặc tính của
CPU. Nội dung thảo luận toàn diện nhiều mặt của tối ưu chương trình, tổng kết
các tư tưởng, nguyên tắc, phương pháp và mẹo thực dụng trong tối ưu, đồng thời
thử đánh giá giá trị áp dụng của chúng trong thi Tin học.
\end{abstract}

\paragraph{Từ khóa.}
Tối ưu; CPU; hợp ngữ; trình biên dịch.

\tableofcontents

\section{Lời nói đầu}

Olympiad in Informatics là cuộc thi về cách viết chương trình máy tính để giải
một bài toán cụ thể. Điểm then chốt là dùng tài nguyên hệ thống hữu hạn, gồm
thời gian CPU và bộ nhớ, để giải các mô hình toán học có quy mô lớn. Khi đã
đặt tính đúng làm tiền đề, ``hiệu quả'' tự nhiên trở thành yếu tố đầu tiên
cần xét. Hiệu quả có hai mặt: thời gian và không gian. Bài viết này tập trung
vào tối ưu hiệu quả thời gian.

\subsection{Thuật toán là yếu tố quyết định}

Nói ngắn gọn, tối ưu thời gian là dùng mọi biện pháp có thể để chương trình
chạy nhanh hơn trong khi vẫn đúng. Công cụ quan trọng nhất vẫn là thuật toán
tốt. Thuật toán là một dãy bước cơ học với ngữ nghĩa xác định, mô tả phương
pháp chung để giải một mô hình toán học. Thuật toán tốt hay xấu trực tiếp
quyết định hiệu quả chạy; khác biệt giữa thuật toán kém và thuật toán tốt
giống như khác biệt giữa đi bộ và đi máy bay, hoàn toàn không cùng bậc.

\subsection{Giới hạn của độ phức tạp thời gian}

Để đo hiệu quả thời gian của thuật toán, ta thường dùng ký hiệu big-O. Khi nói
thời gian chạy bị chặn bởi \(\Oh(f(n))\), nghĩa là nếu thời gian thực tế theo
kích thước dữ liệu \(n\) là \(T(n)\), thì tồn tại các hằng số dương \(n_0,c\)
sao cho với \(n\ge n_0\),
\[
  0\le T(n)\le c f(n).
\]
Cận này giúp ta biết tốc độ tăng của \(T(n)\), từ đó ước lượng thuật toán có
thể chạy trong thời gian hợp lý hay không.

Nhưng công cụ đó rất thô. Hằng số \(c\) trong định nghĩa không quan trọng về
mặt tiệm cận, nhưng trong thực tế lại không thể bỏ qua. Cùng là thuật toán
\(\Oh(n)\), có chương trình chỉ tốn \(2n\) chu kỳ CPU, có chương trình tốn hơn
\(1000n\) chu kỳ. Vì vậy hai cài đặt cùng độ phức tạp vẫn có thể khác nhau rất
xa về thời gian thực.

\subsection{Thuật toán không phải toàn bộ câu chuyện}

Hằng số trong thời gian chạy chịu ảnh hưởng của nhiều yếu tố. Một phần thuộc
tầng logic: thuật toán, cách biểu diễn, cách đơn giản hóa điều kiện. Nhưng còn
một phần nằm dưới tầng logic, khó nhìn thấy hơn nhưng vẫn ảnh hưởng rõ rệt. Ví
dụ, khi viết \texttt{a/=7} trong C, bộ xử lý có thể không dùng phép chia chậm,
mà thay bằng phép nhân và dịch bit.

Vì vậy có thể nhìn hằng số \(c\) trong big-O như \(c=c_1c_2\). \(c_1\) là chi
phí tầng logic, còn \(c_2\) là chi phí tầng thấp của từng câu lệnh. Thời gian
chạy xấp xỉ
\[
  c_2 \cdot [c_1 f(n)].
\]
Phần trong ngoặc phụ thuộc vào thuật toán; \(c_2\) phụ thuộc vào cách chương
trình được triển khai ở tầng dưới. Bài viết này nghiên cứu cách tối ưu những
chi tiết đó.

\subsection{Vì sao cần học tối ưu tầng thấp}

Trong OI, thuật toán là trọng tâm. Tối ưu tầng thấp có vẻ không quan trọng và
không có không gian cải thiện lớn. Tuy vậy khi thuật toán của ta có khiếm
khuyết bẩm sinh, tinh chỉnh tầng thấp có thể giúp đáng kể; phần sau sẽ nêu
vài ví dụ thực tế.

Ngoài thi đấu, hiểu tầng thấp giúp ta hiểu sâu hơn cỗ máy đang dùng. Ngay cả
khi viết bằng ngôn ngữ bậc cao, ta vẫn hình dung được việc gì đang xảy ra bên
dưới, từ đó viết mã có chất lượng tốt hơn.

\subsection{Nền tảng thử nghiệm}

Do điều kiện hạn chế, phần nghiên cứu, lập trình và đo đạc chính được thực
hiện trên một máy ThinkPad X61. Cấu hình liên quan tới tối ưu:
\begin{center}
\begin{tabular}{ll}
CPU & Intel Core 2 Duo T7100 @ 1.80GHz\\
Cache & L1 64KB, L2 2MB\\
Bộ nhớ chính & SODIMM synchronous 667 MHz, 1GB \(\times\) 2\\
Hệ điều hành & Ubuntu 8.10 Intrepid Ibex\\
Nhân & Linux 2.6.27-8-generic\\
Trình biên dịch & gcc 4.3.2\\
Assembler & GNU assembler 2.18.93.20081009\\
Linker & GNU ld 2.18.93.20081009
\end{tabular}
\end{center}

Tác dụng của một kỹ thuật tối ưu phụ thuộc mạnh vào nền tảng, đặc biệt là CPU.
Vì vậy tác giả cũng thử trên nhiều hệ điều hành và nhiều CPU của Intel, AMD.
Giá trị tuyệt đối khác nhau rất nhiều, nhưng hiệu quả tương đối của các hướng
tối ưu cơ bản khá nhất quán, vì CPU hiện đại dùng kiến trúc nội bộ và cơ chế
tối ưu tương tự nhau. Điều cần nắm là nguyên tắc lớn, không phải các mẹo kỳ dị
chỉ hợp một CPU. Riêng nhóm tối ưu \SIMD{} được nhắc ở sau tuy tương thích
trên chip AMD nhưng hiệu quả không nổi bật bằng.

\section{Ví dụ dẫn nhập}

Xét một nhiệm vụ rất đơn giản: viết hàm nhận con trỏ tới đầu mảng số nguyên
dương 32 bit và độ dài mảng, rồi trả về giá trị lớn nhất. Phần đầu hàm trong C
là:
\begin{verbatim}
int get_max(int* a,int l);
\end{verbatim}

\subsection{Cài đặt mộc mạc}

Ở tầng thuật toán không có gì đáng bàn: quét tuyến tính là đủ, và vì mọi phần
tử đều phải được xét để bảo đảm đúng, cận dưới cũng là \(\Oh(n)\).
\begin{verbatim}
// phien ban dau tien
int get_max(int* a,int l){
     int mx=0,i;
     for(i=0;i<l;i++){
          if(a[i]>mx)mx=a[i];
     }
     return mx;
}
\end{verbatim}

Các chương trình trong bài được biên dịch theo cách gần môi trường thi, không
bật tùy chọn tối ưu:
\begin{verbatim}
gcc a.c -o a
\end{verbatim}

Trên mảng \(10\,000\,000\) số nguyên, chạy liên tiếp 100 lần rồi lấy trung
bình, hàm trên tốn \(75\,305\,510\) chu kỳ. Mỗi số mất khoảng 7 đến 8 chu kỳ,
kết quả chưa thỏa đáng.

Hàm đo thời gian dùng lệnh \texttt{rdtsc} của IA-32:
\begin{verbatim}
#define ull unsigned long long
ull get_clock(){
     ull ret;
     __asm__ __volatile__ ("rdtsc\n\t":"=A"(ret):);
     return ret;
}
\end{verbatim}
Phương pháp này khá chính xác nhưng không loại được chu kỳ do hệ điều hành và
tiến trình nền chiếm dụng; có thể kiểm tra chéo bằng lệnh \texttt{time}.

\subsection{Tối ưu lần thứ nhất}

Trong mã đầu, biểu thức \texttt{a[i]} xuất hiện hai lần. Chi phí tính lại địa
chỉ trong vòng lặp nhỏ không phải bằng không, nên ta đổi sang duyệt con trỏ:
\begin{verbatim}
// toi uu lan 1
int get_max(int* a,int l){
     int mx=0,*ed=a+l;
     while(a!=ed){
          if(*a>mx)mx=*a;
          a++;
     }
     return mx;
}
\end{verbatim}

Thời gian trung bình giảm còn \(66\,047\,005\) chu kỳ, tức mỗi phần tử giảm
khoảng một chu kỳ, nhanh hơn bản đầu khoảng \(12\%\). Đây vẫn là tối ưu ở tầng
C và có thể giải thích bằng cách giảm lặp lại phép định địa chỉ.

\subsection{Tối ưu lần thứ hai}

Nếu tiếp tục cùng hướng, ta khó nghĩ ra cải thiện lớn. Điểm mới là tách tiến
trình tìm max thành tám luồng độc lập rồi gộp kết quả:
\begin{verbatim}
// toi uu lan 2
int get_max(int* a,int l){
     assert(l%8==0);
     #define D(x) mx##x=0
     int D(0),D(1),D(2),D(3),
          D(4),D(5),D(6),D(7),*ed=a+l;
     #define CMP(x) if(*(a+x)>mx##x)mx##x=*(a+x);
     while(a!=ed){
          CMP(0);CMP(1);
          CMP(2);CMP(3);
          CMP(4);CMP(5);
          CMP(6);CMP(7);
          a+=8;
     }
     #define CC(x1,x2) if(mx##x1>mx##x2)mx##x2=mx##x1;
     CC(1,0);CC(3,2);
     CC(5,4);CC(7,6);
     CC(2,0);CC(6,4);
     CC(4,0);
     return mx0;
}
\end{verbatim}

Mã dùng nhiều macro và chỉ xử lý trực tiếp trường hợp \(l\) chia hết cho 8,
nhưng ý tưởng rất đơn giản. Kết quả đo là \(34\,818\,706\) chu kỳ, nhanh hơn
bản đầu \(54\%\). Nếu chỉ nhìn tầng C, khó giải thích mức cải thiện này: số
lần tăng con trỏ giảm xuống \(1/8\), nhưng thao tác đó nhỏ hơn nhiều so với
phần còn lại của vòng lặp.

Điểm mấu chốt là mã sau khi biên dịch. Ở bản đầu, mỗi lần cập nhật
\texttt{mx} phụ thuộc vào kết quả lần trước, nên các lệnh liên quan buộc phải
đi nối tiếp. Khi tách thành \texttt{mx0}, \texttt{mx1}, \ldots, các chuỗi lệnh
không còn phụ thuộc nhau, CPU có thêm cơ hội chạy song song bằng cơ chế nội
bộ như pipeline, dự đoán nhánh, đổi tên thanh ghi, cache và thực thi ngoài
thứ tự. Chi tiết sẽ được giải thích ở các phần sau.

\subsection{Tối ưu lần thứ ba: hợp ngữ nội tuyến}

Khi đã nhìn tới tầng hợp ngữ, ta có thể viết trực tiếp đoạn nóng bằng hợp ngữ.
Trong CPU, các khái niệm như thanh ghi, cờ trạng thái và lệnh là trung tâm,
nhưng bị ngôn ngữ bậc cao che đi. C cũng dựa nhiều vào biến trong bộ nhớ, mà
đọc ghi bộ nhớ là một trong các thao tác chậm. Hợp ngữ cho phép thao tác trực
tiếp hơn và dùng các lệnh đặc biệt. Đổi lại, chi phí lập trình rất cao, khó
bảo trì và phụ thuộc nền tảng. Cách thực dụng trong thi là dùng hợp ngữ nội
tuyến cho đoạn nghẽn cổ chai, còn phần chính vẫn viết bằng C.

Các đoạn sau dùng hợp ngữ nội tuyến của gcc, cú pháp AT\&T:
\begin{verbatim}
// toi uu lan 3
int get_max(int* a,int l){
     int ret;
     __asm__ __volatile__ (
           "movl      $0, %%eax\n\t"
           ".p2align 4,,15\n"
           "LP1:\n\t"
           "cmpl      -4(%1,%2,4), %%eax\n\t"
           "jge ED\n\t"
           "movl      -4(%1,%2,4), %%eax\n"
           "ED:\n\t"
           //"loop LP1\n\t"
           "decl      %2\n\t"
           "jnz LP1\n\t"
           "movl      %%eax, %0\n\t"
           :"=m"(ret)
           :"r"(a),"c"(l)
           :"%eax");
     return ret;
}
\end{verbatim}

Thời gian trung bình là \(21\,322\,853\) chu kỳ, chỉ khoảng 2 chu kỳ cho một
phần tử, nhanh hơn bản đầu \(72\%\). Vòng lặp lõi vẫn có 5 lệnh, trong đó có
hai nhánh điều kiện và hai lần truy cập bộ nhớ bằng định địa chỉ SIB phức tạp.
Kết quả tốt chủ yếu nhờ cache dữ liệu, dự đoán nhánh và thực thi ngoài thứ tự.

\paragraph{\RISC{} và \CISC.}
Trong mã trên có một lệnh \texttt{loop} bị chú thích. Trực giác cũ có thể nghĩ
\texttt{loop} ngắn hơn và phức hợp hơn nên nhanh hơn. Thực tế thay
\texttt{decl}/\texttt{jnz} bằng \texttt{loop} làm thời gian tăng lên
\(56\,457\,348\) chu kỳ.

Lý do nằm ở lịch sử \CISC{} và \RISC. x86 là kiến trúc \CISC{} với tập lệnh
lớn và nhiều lệnh phức tạp. Trên CPU hiện đại, nhiều lệnh x86 được dịch bên
trong thành các vi lệnh giống \RISC. Các nhà sản xuất tối ưu rất mạnh cho
nhóm lệnh đơn giản phù hợp pipeline, trong khi một số lệnh phức tạp cũ như
\texttt{loop} không còn có ưu thế. Vì vậy khi tối ưu hợp ngữ hiện nay, thường
nên dùng tập con lệnh cơ bản đã được tối ưu tốt. Ngoại lệ quan trọng là các
tập lệnh \SIMD{} được đầu tư riêng.

\subsection{Tối ưu lần thứ tư: giảm phụ thuộc}

Ta thử lặp lại ý tưởng tách chuỗi phụ thuộc trong hợp ngữ:
\begin{verbatim}
// toi uu lan 4
int get_max(int* a,int l){
     assert(l%2==0);
     int ret;
     __asm__ __volatile__ (
           "movl      $0, %%eax\n\t"
           "movl      $0, %%edx\n\t"
           ".p2align 4,,15\n"
           "LP2:\n\t"
           "cmpl      (%1),   %%eax\n\t"
           "jge ED2\n\t"
           "movl      (%1),   %%eax\n"
           "ED2:\n\t"
           "cmpl      4(%1), %%edx\n\t"
           "jge ED3\n\t"
           "movl      4(%1), %%edx\n"
           "ED3:\n\t"
           "addl      $8, %1\n\t"
           "subl      $2, %2\n\t"
           "jnz LP2\n\t"
           "cmpl      %%edx, %%eax\n\t"
           "cmovll %%edx, %%eax\n\t"
           "movl      %%eax, %0\n\t"
           :"=m"(ret)
           :"r"(a),"r"(l)
           :"%eax","%edx");
     return ret;
}
\end{verbatim}

Kết quả là \(17\,447\,544\) chu kỳ, có cải thiện nhưng không thần kỳ như bản
C tám luồng. Khi thử mở rộng bốn luồng và dùng hết thanh ghi tổng quát, tốc độ
lại giảm. Nguyên nhân là bản hợp ngữ đã rất gọn; khi lệnh tải đầu vòng lặp bị
cache miss, phía sau không còn nhiều lệnh độc lập để CPU kéo lên chạy trước.
Thêm quá nhiều nhánh và dùng hết thanh ghi còn gây bất lợi cho RAT và dự đoán
nhánh. Tối ưu không thể dựa vào tưởng tượng.

\subsection{Tối ưu lần thứ năm: \SIMD}

\SIMD{} là ``single instruction, multiple data'': một lệnh xử lý đồng thời
nhiều phần tử dữ liệu. Nó khác với lệnh phức hợp kiểu ``một lệnh làm nhiều
việc''; ở đây cùng một thao tác được áp vào nhiều phần tử đã được đóng gói
trong thanh ghi rộng. Intel bắt đầu từ MMX, rồi SSE với thanh ghi 128 bit, sau
đó SSE2, SSE3, SSE4; AMD cũng tương thích rộng rãi.

\begin{verbatim}
// toi uu lan 5
int get_max(int* a,int l){
     assert(l%4==0);
     assert(sse2);
     int ret,tmp[4];
     __asm__ __volatile__ (
           "\txorps %%xmm0,         %%xmm0\n"
           "LP3:\n"
           "\tmovdqa        %%xmm0,     %%xmm1\n"
           "\tpcmpgtd       (%1),   %%xmm1\n"
           "\tandps %%xmm1,         %%xmm0\n"
           "\tandnps (%1),       %%xmm1\n"
           "\torps          %%xmm1,     %%xmm0\n"
           "\taddl $16, %1\n"
           "\tsubl $4, %2\n"
           "\tjnz     LP3\n"
           "\tmovdqu        %%xmm0,     (%3)\n"
           "\tmovl (%3),         %%eax\n"
           "\tcmpl 4(%3), %%eax\n"
           "\tcmovll 4(%3), %%eax\n"
           "\tcmpl 8(%3), %%eax\n"
           "\tcmovll 8(%3), %%eax\n"
           "\tcmpl 12(%3), %%eax\n"
           "\tcmovll 12(%3), %%eax\n"
           "\tmovl %%eax, %0\n"
           :"=m"(ret)
           :"r"(a),"r"(l),"r"(tmp)
           :"%eax");
     return ret;
}
\end{verbatim}

Bốn giá trị lớn nhất cục bộ được đóng trong \texttt{xmm0}. Để lấy max giữa
bốn số liên tiếp trong bộ nhớ và bốn số trong \texttt{xmm0}, mã dùng so sánh
SSE tạo mặt nạ rồi ghép bằng phép bit. Dù cần tới 5 lệnh SSE, thời gian giảm
còn \(15\,898\,751\) chu kỳ. Bản này chỉ tốn khoảng \(1/5\) thời gian bản đầu,
tức xấp xỉ 1.5 chu kỳ cho mỗi phần tử.

\subsection{Tối ưu lần thứ sáu: SSE4}

Trên các CPU Penryn 45nm của Intel, SSE4 có lệnh \texttt{pmaxsd}, có thể thay
trực tiếp khối so sánh và phép bit:
\begin{verbatim}
// toi uu lan 6
int get_max(int* a,int l){
     assert(l%4==0);
     assert(sse4);
     int ret,tmp[4];
     __asm__ __volatile__ (
           "\txorps %%xmm0,         %%xmm0\n"
           "LP4:\n"
           "\tpmaxsd        (%1),   %%xmm0\n"
           "\taddl $16, %1\n"
           "\tsubl $4, %2\n"
           "\tjnz     LP4\n"
           "\tmovdqu        %%xmm0,     (%3)\n"
           "\tmovl (%3),         %%eax\n"
           "\tcmpl 4(%3), %%eax\n"
           "\tcmovll 4(%3), %%eax\n"
           "\tcmpl 8(%3), %%eax\n"
           "\tcmovll 8(%3), %%eax\n"
           "\tcmpl 12(%3), %%eax\n"
           "\tcmovll 12(%3), %%eax\n"
           "\tmovl %%eax, %0\n"
           :"=m"(ret)
           :"r"(a),"r"(l),"r"(tmp)
           :"%eax");
     return ret;
}
\end{verbatim}
Máy thử nghiệm của tác giả chưa chạy được SSE4 nên không có số đo, nhưng tác
giả ước đoán nó có thể tiến gần mức cực hạn một chu kỳ cho mỗi phần tử.

\section{Hiệu quả chạy của các lệnh CPU}

Trên kiến trúc \CISC, rất khó nói một lệnh cụ thể luôn tốn bao nhiêu chu kỳ:
vi lệnh của nó có thể bị tách ra và chạy xen với lệnh khác; dữ liệu có ở
cache hay không, có căn chỉnh hay không cũng ảnh hưởng lớn. Tuy vậy vẫn cần
biết đại khái lệnh nào nhanh chậm. Một lệnh chia gần như luôn chậm hơn cộng.

Tác giả đo bằng cách đặt lệnh cần kiểm tra trong một ngữ cảnh có ý nghĩa, chèn
\texttt{cpuid} trước và sau để cô lập, lặp \(1\,000\,000\) lần, rồi trừ chi
phí vòng lặp không có lệnh đó. Kết quả trên máy thử nghiệm:
\begin{center}
\begin{tabular}{lll}
\textbf{Lệnh} & \textbf{Toán hạng} & \textbf{Chu kỳ trung bình}\\
\texttt{add} & \texttt{r32/r32} & 1\\
\texttt{shr} & \texttt{imm8/r32} & 1\\
\texttt{bsr} & \texttt{r32/r32} & 1\\
\texttt{mul} & \texttt{r32} & 2\\
\texttt{div} & \texttt{r32} & 21\\
\texttt{fadd} & \texttt{r} & 3\\
\texttt{fdiv} & \texttt{r} & 35\\
\texttt{fsqrt} & n/a & 60\\
\texttt{padd} & \texttt{xmm/xmm} & 1
\end{tabular}
\end{center}

\section{Tối ưu phép toán số học}

\subsection{Phép chia}

Trên CPU hiện đại, phép nhân đã được tối ưu rất nhanh, nhưng phép chia vẫn
phức tạp và tốn hàng chục chu kỳ. Nếu 20 lệnh dịch bit có thể thay một lệnh
chia, việc đó thường đáng làm. Tuy nhiên không tồn tại thuật toán chia tổng
quát nhanh hơn phần cứng CPU; nếu có, nhà sản xuất đã đưa nó vào mạch chia.
Không gian tối ưu chỉ xuất hiện trong các trường hợp đặc biệt.

\subsubsection{Loại bỏ phép chia}

Cách hiệu quả nhất là tránh chia. Ví dụ điều kiện \texttt{if(a/b<c)} đôi khi
có thể đổi thành \texttt{if(a<b*c)}, nhưng phải xét tràn số và dấu của
\texttt{b}. Trong thiết kế thuật toán cũng có những biến đổi tương tự, nhưng
cần rất cẩn thận.

Trong các bài đếm modulo \(M\), lạm dụng \texttt{\%M} có thể rất đắt vì phép
lấy dư ở tầng dưới dùng cùng lệnh với phép chia. Chẳng hạn để chuẩn hóa một
số có thể âm, cách thường viết \texttt{(a\%M+M)\%M} dùng hai phép chia. Có thể
viết:
\begin{verbatim}
inline int mod(int a){a%=M;if(a<0)a+=M;return a;};
\end{verbatim}

Trong nhân ma trận cho truy hồi tuyến tính, nếu lấy modulo ở mọi bước trong
vòng lặp trong cùng, chương trình sẽ chậm đáng kể. Thường có thể cộng dồn tạm
thời trong phạm vi an toàn rồi mới lấy modulo sau khi tính xong một phần tử.

\subsubsection{Giảm chia số thực}

Với số thực, biến đổi linh hoạt hơn vì quy luật đại số quen thuộc hơn. Nếu
cần chia nhiều lần cho \(b\), nên tính trước \(1/b\) rồi nhân. Với
\[
  a/b+c/d,
\]
có thể đổi thành
\[
  (a d+c b)/(b d),
\]
tăng 3 phép nhân nhưng giảm 1 phép chia; trên máy tác giả cách này nhanh hơn
hơn hai lần khi lặp \(10\,000\,000\) lần.

Ngay cả ba phép chia độc lập:
\begin{verbatim}
x=a/b;
y=c/d;
z=e/f;
\end{verbatim}
có thể đổi thành:
\begin{verbatim}
t=1/(b*d*f);
x=a*d*f*t;
y=c*b*f*t;
z=e*b*d*t;
\end{verbatim}
Mã thứ hai chỉ tốn khoảng \(1/3\) thời gian mã đầu trong thử nghiệm. Dĩ nhiên
không nên lạm dụng: quá nhiều lệnh đơn giản thay một lệnh phức tạp có thể làm
mã khó đọc, khó kiểm chứng và ảnh hưởng độ chính xác.

\subsubsection{Các phép chia đặc biệt}

Nếu toán hạng chắc chắn không âm, nên chuyển sang kiểu không dấu trước khi
chia:
\begin{verbatim}
a=(unsigned int)b/123;
\end{verbatim}
Chia không dấu thường nhanh hơn chia có dấu.

\paragraph{Chia cho lũy thừa của 2.}
Chia cho \(2^k\) tương đương dịch phải \(k\) bit. Nếu modulo cũng là
\(2^k\), có thể để số tràn tự nhiên và chỉ giữ \(k\) bit thấp.

Đôi khi cần tìm \(k\) khi biết \(b=2^k\). Một mẹo C dùng bảng băm không va
chạm:
\begin{verbatim}
char tb[32]={31,0,27,1,28,18,23,2,
              29,21,19,12,24,9,14,3,
              30,26,17,22,20,11,8,13,
              25,16,10,7,15,6,5,4};
// x can la unsigned int
#define LOG(x) tb[((x)*263572066)>>27]
\end{verbatim}
Ý tưởng là chọn một số \(T\) sao cho khi dịch trái \(T\) lần lượt 0 đến 31
bit, 5 bit cao chạy qua đủ các giá trị 0 đến 31.

Nhưng IA-32 đã có lệnh quét bit. \texttt{bsf} quét từ thấp lên cao, còn
\texttt{bsr} từ cao xuống thấp:
\begin{verbatim}
inline unsigned int LOG2(unsigned int x){
     unsigned int ret;
     __asm__ __volatile__ ("bsrl %1, %%eax":"=a"(ret):"m"(x));
     return ret;
}
\end{verbatim}
\texttt{LOG2} không chỉ xử lý lũy thừa của 2 mà trả về \(\lfloor\log_2 x\rfloor\).
Nó hữu ích cho bảng sparse table, LCA và nhiều thuật toán tìm kiếm theo bit.
\texttt{bsf} hữu ích để đếm số bit 0 ở cuối, ví dụ trong Miller-Rabin hoặc
GCD nhị phân.

\paragraph{Chia cho hằng số.}
Với số thực, chia cho hằng số chỉ cần nhân với nghịch đảo. Với số nguyên,
không thể làm trực tiếp như vậy, nhưng trình biên dịch có kỹ thuật riêng. Khi
biên dịch:
\begin{verbatim}
#include<stdio.h>
int main(){
     unsigned int a;
     scanf("%u",&a);
     a/=7;
     printf("%u",a);
     return 0;
}
\end{verbatim}
gcc \texttt{-O2} sinh đoạn tương đương:
\begin{verbatim}
movl      -8(%ebp), %ecx
movl      $613566757, %edx
movl      %ecx, %eax
mull %edx
subl %edx, %ecx
shrl %ecx
addl %ecx, %edx
shrl $2, %edx
movl      %edx, -8(%ebp)
\end{verbatim}

Con số \(613566757\) thỏa \(613566757\cdot 7=2^{32}+3\). Bằng cách nhân vào
định nghĩa phép chia có dư \(7a+q=c\), các phần ``lộn xộn'' chứa \(q\) nằm ở
bit thấp và bị loại sau khi dịch phải đủ nhiều, còn thương xuất hiện ở bit
cao. gcc dùng vài biến thể tùy hằng số. Dù gcc đã làm việc này ngay cả khi
không bật tối ưu mạnh, hiểu nguyên lý cho thấy có thể mở rộng sang trường hợp
mẫu số được biết muộn nhưng dùng lặp lại, bằng cách tự sinh mã máy; tuy vậy
cách đó quá đắt cho thi đấu.

\paragraph{Dùng trừ thay chia?}
Một số mã GCD cố thử trừ \(b\) khỏi \(a\) vài lần trước khi dùng phép chia.
Tác giả thử nghiệm và thấy chậm hơn, kể cả khi viết hợp ngữ. Lý do là tốc độ
lệnh chia không cố định; khi \(a\) và \(b\) gần nhau, nó có thể khá nhanh.
Trong khi đó mô phỏng bằng trừ tạo nhiều nhánh điều kiện, làm CPU khó chạy
trơn tru. Đây là ví dụ điển hình rằng tối ưu phải đo, không suy diễn.

\subsection{Phép nhân}

So với chia, phép nhân ít chỗ để nghiên cứu hơn. Trên CPU hiện đại,
\texttt{mul} đã đủ nhanh; nhiều tối ưu cổ điển như viết \(a\cdot 10\) thành
\texttt{(a<<3)+(a<<1)} có thể chậm hơn do cần hai dịch, một cộng và vài
\texttt{mov} ẩn.

\subsubsection{Dùng \texttt{lea} cho nhân hằng số}

gcc nhân 10 có thể sinh:
\begin{verbatim}
lea (%eax,%eax,4),%eax
add %eax,%eax
\end{verbatim}
\texttt{lea offset(reg1,reg2,c),dest} tính \(reg1+reg2\cdot c+offset\), với
\(c\in\{1,2,4,8\}\). Vì vậy dòng đầu nhân \texttt{eax} với 5, dòng sau nhân
đôi thành 10. Một số ví dụ do gcc sinh:
\begin{verbatim}
nhan 2:  addl %eax,%eax
nhan 9:  leal (%eax,%eax,8),%eax
nhan 13: leal (%eax,%eax,2),%edx
         leal (%eax,%edx,4),%edx
nhan 17: movl %eax,%edx
         sall 4,%edx
         addl %edx,%eax
nhan 68: movl %eax,%edx
         sall 6,%edx
         leal (%edx,%eax,4),%eax
\end{verbatim}
Không phải hằng số nào cũng có biểu diễn ngắn; với \texttt{a*14}, gcc dùng
\texttt{imul}.

\subsubsection{Tối ưu \texttt{mul\_mod}}

Một sự thật dễ bị che đi ở C là lệnh \texttt{mul} không ``tràn'' theo nghĩa
mất kết quả: nhân hai số 32 bit cho kết quả 64 bit trong \texttt{edx:eax}.
Lệnh \texttt{div} cũng lấy bị chia từ \texttt{edx:eax}. Vì vậy có thể viết:
\begin{verbatim}
#define MO 123456789
inline int mul_mod(int a,int b){
     int ret;
     __asm__ __volatile__ ("\tmull %%ebx\n\tdivl %%ecx\n"
                 :"=d"(ret):"a"(a),"b"(b),"c"(MO));
     return ret;
}
\end{verbatim}
Trên máy tác giả, hàm này nhanh hơn hơn hai lần so với:
\begin{verbatim}
inline int mul_mod_slow(int a,int b){
     return (long long)a*b%MO;
}
\end{verbatim}

Khi modulo là số 64 bit, có một mẹo dùng số thực để xấp xỉ thương:
\begin{verbatim}
typedef long long ll;
#define MOL 123456789012345LL
inline ll mul_mod_ll(ll a,ll b){
     ll d=(ll)floor(a*(double)b/MOL+0.5);
     ll ret=a*b-d*MOL;
     if(ret<0)ret+=MOL;
     return ret;
}
\end{verbatim}
Hai phép nhân trong \texttt{a*b-d*MOL} có thể tràn, nhưng hiệu của chúng được
dự đoán nằm trong 64 bit; phần tràn chỉ lệch 0 hoặc 1 lần modulo, rồi được
sửa bằng kiểm tra dấu. Không nên dùng với modulo quá lớn vì sai số số thực.

\subsection{Số nguyên lớn}

Theo tác giả, mã số nguyên lớn vốn rất hợp với hợp ngữ. IA-32 có nhiều lệnh
gần như được sinh ra cho số nguyên lớn; ở tầng hợp ngữ có thể lưu số theo cơ
số \(2^{32}\) để tận dụng lệnh và đạt tốc độ cao. Tác giả từng viết một thư
viện nội tuyến nhanh hơn cách truyền thống vài lần.

\section{Đặc tính tối ưu của CPU}

CPU không chỉ đơn giản thực hiện từng lệnh theo thứ tự. Khi tần số xung bị
giới hạn bởi công nghệ chế tạo và công suất, hiệu năng vẫn tăng nhờ các cơ
chế như nạp trước lệnh, dự đoán nhánh, thực thi ngoài thứ tự và cache. Cùng
một chức năng nhưng cách viết hơi khác có thể khiến CPU tối ưu tốt hoặc kém
rất xa. Phần này chỉ giữ các điểm có ý nghĩa ở tầng ngôn ngữ bậc cao.

\subsection{Cache}

\subsubsection{Truy cập bộ nhớ chậm}

Biến trong ngôn ngữ bậc cao thường nằm trong bộ nhớ, còn tốc độ bộ nhớ kém xa
tốc độ CPU. Khi cần dữ liệu, CPU phải gửi yêu cầu qua bus, đợi bộ nhớ chuẩn
bị rồi lấy về, có thể tốn hàng chục hoặc hàng trăm chu kỳ. CPU giải quyết
bằng hai hướng: trong lúc đợi thì làm việc khác, và dùng cache.

Cache dựa trên tính cục bộ: chương trình thường thực thi mã gần nhau và truy
cập dữ liệu gần nhau; một vị trí vừa truy cập có xác suất cao sẽ được dùng
lại sớm. Máy thử nghiệm có L1 64KB, L2 2MB; truy cập L1 gần như truy cập
thanh ghi.

Khi tải một biến vào cache, CPU không tải từng byte riêng lẻ mà tải cả một
dòng cache. Trên máy tác giả, mỗi dòng là 64 byte. L1 có 1024 dòng nhưng một
địa chỉ chỉ có thể vào một nhóm nhất định; mỗi nhóm có 8 dòng. Nếu nhóm đầy,
phải đẩy một dòng ra. Truy cập đã có trong cache là hit; chưa có là cache
miss, rất đắt.

\subsubsection{Thí nghiệm quan sát cache}

\begin{verbatim}
int sum=0;
for(i=0;i<N;i++)sum+=a[i];
\end{verbatim}
Với \(N=10\,000\,000\), biên dịch \texttt{-O2}, chạy mười lần:
\begin{center}
\begin{tabular}{lr}
Lần & Chu kỳ\\
0 & 167150043\\
1 & 24793497\\
2 & 27947781\\
3 & 25602750\\
4 & 27184617\\
5 & 18583821\\
6 & 15880428\\
7 & 15828174\\
8 & 15944328\\
9 & 16141635
\end{tabular}
\end{center}
Lần đầu chậm hơn khoảng một bậc vì mảng chưa được nạp vào cache.

\subsubsection{Viết mã thân thiện với cache}

\paragraph{Ưu tiên biến cục bộ.}
Dữ liệu trên stack được truy cập rất thường xuyên và tập trung; đỉnh stack
thường nằm trong L1. Biến toàn cục cần địa chỉ 32 bit, làm mã lệnh dài hơn và
tỷ lệ cache hit thường kém hơn.

\paragraph{Dùng \texttt{struct} cho dữ liệu liên quan.}
Hai mảng tách rời:
\begin{verbatim}
int x[MAXN],y[MAXN];
\end{verbatim}
thường kém hơn:
\begin{verbatim}
struct point{int x,y;};
point p[MAXN];
\end{verbatim}
vì khi dùng \(x\) của một điểm, ta thường sắp dùng \(y\) của cùng điểm; đặt
gần nhau giúp cùng vào một dòng cache.

\paragraph{Căn chỉnh dữ liệu.}
Nếu một \texttt{struct point} bị vắt qua ranh giới dòng cache, mỗi lần truy
cập có thể phải tải hai dòng. gcc cho phép yêu cầu căn chỉnh:
\begin{verbatim}
int array[1024] __attribute__((aligned(64)));
\end{verbatim}

\paragraph{Dùng \texttt{union} khi phù hợp.}
Nếu hai dữ liệu không bao giờ được dùng đồng thời, cho chúng chồng lên cùng
vị trí bằng \texttt{union} vừa tiết kiệm bộ nhớ vừa tận dụng cache khi chuyển
từ dữ liệu này sang dữ liệu kia.

\paragraph{Tránh cấp phát động dày đặc.}
Cấp phát và giải phóng động phải quản lý danh sách khối nhớ, các khối sinh ra
có thể rời rạc và bất lợi cho cache. Với cấu trúc dữ liệu động trong thi, nên
mở trước một pool đủ lớn rồi tự cấp phát node. Node nên là \texttt{struct}
thay vì nhiều mảng rời.

Không phải mọi mảng đều nên đặt trên stack. Mảng rất lớn dễ tràn stack và làm
con trỏ stack dao động mạnh, phá lợi thế cục bộ. Tính năng mảng cục bộ có
kích thước runtime của gcc cũng không được khuyến nghị trong ngữ cảnh này.

\subsubsection{Cache mã lệnh}

Mã chương trình cũng nằm trong bộ nhớ. CPU có cache lệnh tương tự cache dữ
liệu. Ở tầng C, hai nguyên tắc chính là đặt các hàm liên quan gần nhau và giữ
vòng lặp nóng gọn.

Ví dụ:
\begin{verbatim}
for(int i=0;i<1000000;i++){
     proc1();
     proc2();
}
\end{verbatim}
nếu \texttt{proc1} và \texttt{proc2} phức tạp và ít liên quan, có thể tách:
\begin{verbatim}
for(int i=0;i<1000000;i++)proc1();
for(int i=0;i<1000000;i++)proc2();
\end{verbatim}
Dù lặp bộ đếm hai lần, mỗi vòng lặp nhỏ hơn, mã và dữ liệu liên quan có cơ
hội ở cache tốt hơn; dự đoán nhánh cũng dễ hơn. Nguyên tắc tổng quát: thứ có
liên quan nên đặt gần, thứ không liên quan nên tách ra.

\subsection{Dự đoán nhánh}

Điều kiện trong C được biên dịch thành các lệnh nhảy có điều kiện. Nhánh có
điều kiện gây khó cho cache mã, trace cache và đặc biệt là pipeline. Một lệnh
hiện đại đi qua nhiều bước: nạp trước, giải mã, đổi tên thanh ghi, thực thi,
retire. Nếu tới một nhánh mà CPU không biết đường nào đúng, nó có thể dừng
pipeline chờ kết quả, rất đắt. Do đó CPU dự đoán hướng nhánh và chạy trước.
Nếu dự đoán sai, CPU phải quay lại và làm lại, có thể mất hàng chục chu kỳ.

Thí nghiệm:
\begin{verbatim}
int main(){
     for(int i=0;i<N;i++)a[i]=i%2;
     ull start=get_clock();
     for(int i=0;i<N;i++){
           if(a[i])a[i]=123123123;
           else a[i]=1234123;
     }
     printf("%llu\n",get_clock()-start);
     return 0;
}
\end{verbatim}
Kết quả là \(32\,889\,996\) chu kỳ. Nếu đổi dòng khởi tạo thành:
\begin{verbatim}
for(int i=0;i<N;i++)a[i]=rand()%2;
\end{verbatim}
kết quả tăng lên \(102\,583\,557\). Mã đo không đổi; khác biệt chỉ là mẫu dữ
liệu. Chuỗi 0,1 xen kẽ có quy luật nên CPU học được, còn dữ liệu ngẫu nhiên
thì không.

Dự đoán nhánh thường tốt khi nhánh luôn đi cùng một hướng, lặp một mẫu đơn
giản, hoặc tương quan với một nhánh trước đó. Vì vậy nhánh do so sánh bộ đếm
vòng lặp thường an toàn: hầu như luôn tiếp tục, chỉ sai ở lần thoát.

\subsubsection{Loại bỏ nhánh}

Một số ví dụ hợp ngữ:
\begin{verbatim}
// tri tuyet doi cua eax
cltd
xorl %edx, %eax
subl %edx, %eax
\end{verbatim}
Nếu \texttt{eax} không âm, \texttt{edx=0}; nếu âm, \texttt{edx=-1}, phép xor
rồi trừ biến nó thành bù hai của số âm.

Gán \texttt{ebx=eax} nếu \texttt{ebx>eax}:
\begin{verbatim}
subl %ebx,    %eax
sbbl %edx,    %edx
andl %eax,    %edx
addl %edx,    %ebx
\end{verbatim}

Hàm so sánh dấu:
\begin{verbatim}
int cmp(int a){return (a>>31)+(-a>>31&1);}
\end{verbatim}

Trong ngôn ngữ bậc cao, mỗi hạng tử trong biểu thức logic có thể sinh nhánh
riêng, nên điều kiện trong \texttt{if} nên được viết gọn và rõ.

\subsection{Thực thi ngoài thứ tự}

Thực thi ngoài thứ tự (\textit{out-of-order execution}) cố sắp lại các vi
lệnh để tận dụng các đơn vị thực thi song song. Hai kỹ thuật nền tảng là vi
lệnh và đổi tên thanh ghi.

\paragraph{Vi lệnh.}
Lệnh x86 đi vào CPU được giải mã thành các thao tác nhỏ kiểu \RISC. Ví dụ:
\begin{verbatim}
addl (%esi),   %eax
\end{verbatim}
có thể tách thành một vi lệnh tải bộ nhớ và một vi lệnh cộng. Còn:
\begin{verbatim}
addl %eax, (%esi)
\end{verbatim}
cần tải bộ nhớ, cộng, rồi ghi lại bộ nhớ. Khi đã tách thành vi lệnh, CPU có
thể đưa chúng vào các cổng khác nhau: tải bộ nhớ, ghi bộ nhớ, số nguyên, số
thực. Nếu một vi lệnh đang đợi cache miss, các vi lệnh độc lập phía sau có
thể chạy trước.

\paragraph{Đổi tên thanh ghi.}
Các thanh ghi logic như \texttt{eax} chỉ là tên bên ngoài. Bên trong CPU có
nhiều thanh ghi vật lý hơn; RAT ánh xạ thanh ghi logic sang thanh ghi vật lý.
Hai đoạn mã độc lập cùng dùng \texttt{eax} có thể được đổi thành hai thanh ghi
vật lý khác nhau, giúp chúng chạy ngoài thứ tự.

Quy trình chung là: đơn vị nạp trước đọc lệnh, giải mã và đổi tên thanh ghi,
đưa vào reorder buffer để sắp lại, gửi tới đơn vị thực thi, giữ kết quả, rồi
retirement unit trả kết quả theo đúng thứ tự chương trình. Nếu bước nạp trước
sai vì dự đoán nhánh sai, phần sau phải hủy hết, cho thấy nhánh nguy hiểm thế
nào.

\subsubsection{Giảm phụ thuộc giữa câu lệnh}

Ở C, ta không kiểm soát vi lệnh, nhưng có một nguyên tắc vẫn dùng được: giảm
phụ thuộc giữa các câu lệnh. Ví dụ:
\begin{verbatim}
double sum=0;
for(int i=0;i<N;i++)sum+=a[i];
\end{verbatim}
với \(N=10\,000\,000\) tốn \(88\,532\,892\) chu kỳ. Đổi thành bốn luồng cộng:
\begin{verbatim}
double s0=0,s1=0,s2=0,s3=0;
for(int i=0;i<N;i+=4){
     s0+=a[i];s1+=a[i+1];
     s2+=a[i+2];s3+=a[i+3];
}
double sum=(s0+s1)+(s2+s3);
\end{verbatim}
thời gian giảm còn \(43\,307\,199\) chu kỳ. Chuỗi phụ thuộc trên \texttt{sum}
đã được tách thành bốn chuỗi độc lập, giúp CPU chạy song song tốt hơn.

Một nguyên tắc phụ: trên máy 32 bit, nếu không thiếu bộ nhớ, nên dùng kiểu
nguyên 32 bit. Các thanh ghi 16 bit như \texttt{ax} chỉ là phần của thanh ghi
32 bit, khó đổi tên hơn và mã lệnh thường dài hơn.

\section{Kỹ thuật bit}

Mọi dữ liệu trong máy tính đều được lưu nhị phân. Các phép logic bit như
\texttt{and}, \texttt{or}, \texttt{not}, \texttt{xor}, dịch bit và quét bit
gần với cổng logic phần cứng nhất, thường chạy trong một chu kỳ. Kết hợp khéo
các phép này có thể thay thao tác chậm hơn.

\subsection{Nén bit}

Giả sử \texttt{a} là biến 32 bit chứa 32 cờ boolean:
\begin{center}
\begin{tabular}{ll}
Đọc bit \(k\) & \texttt{a>>k\&1}\\
Đọc bit \(k\) rồi đảo & \texttt{\textasciitilde{}a>>k\&1}\\
Xóa bit \(k\) & \texttt{a\&=\textasciitilde{}(1<<k)}\\
Bật bit \(k\) & \texttt{a|=1<<k}\\
Đảo bit \(k\) & \texttt{a\string^=1<<k}\\
Đảo đoạn \(k_1..k_2\) & \texttt{a\string^=((1<<(k2-k1+1))-1)<<k2}\\
Đúng một bit bật & \texttt{!(x\&(x-1))\&\&x}\\
Có hai bit bật kề nhau & \texttt{x>>1\&x}\\
Có ba bit bật kề nhau & \texttt{x>>1\&x>>2\&x}
\end{tabular}
\end{center}

\subsection{Thống kê trong một gói bit}

Đếm từng bit sẽ mất lợi thế của nén. Có thể dùng bảng nhỏ, chẳng hạn
\(2^{16}\), hoặc xử lý bằng phép bit theo kiểu chia đôi.

Đếm parity của số bit bật:
\begin{verbatim}
x^=x>>1;x^=x>>2;
x^=x>>4;x^=x>>8;
x^=x>>16;
\end{verbatim}
Sau đó bit \(i\) của kết quả là parity của đoạn từ \(i\) tới bit cao nhất.
Parity của đoạn \(k_1..k_2\) là:
\begin{verbatim}
(x>>k1^x>>(k2+1))&1
\end{verbatim}

Đếm số bit bật:
\begin{verbatim}
int count(unsigned int x)
{
     x=(x&0x55555555)+(x>>1&0x55555555);
     x=(x&0x33333333)+(x>>2&0x33333333);
     x=(x&0x0F0F0F0F)+(x>>4&0x0F0F0F0F);
     x=(x&0x00FF00FF)+(x>>8&0x00FF00FF);
     x=(x&0x0000FFFF)+(x>>16&0x0000FFFF);
     return x;
}
\end{verbatim}
Đối số nên là \texttt{unsigned int}. Với kiểu có dấu, dịch phải có thể là
dịch số học (\texttt{sar}) và mở rộng kiểu có thể điền bit dấu, rất dễ gây
lỗi khó tìm.

Đảo thứ tự bit:
\begin{verbatim}
unsigned int rev(unsigned int x){
     x=(x&0x55555555)<<1|(x>>1&0x55555555);
     x=(x&0x33333333)<<2|(x>>2&0x33333333);
     x=(x&0x0F0F0F0F)<<4|(x>>4&0x0F0F0F0F);
     x=(x&0x00FF00FF)<<8|(x>>8&0x00FF00FF);
     x=(x&0x0000FFFF)<<16|(x>>16&0x0000FFFF);
     return x;
}
\end{verbatim}

\subsection{Loại bỏ nhánh bằng C}

\begin{verbatim}
int abs(int x){
     int y=x>>31;
     return (x+y)^y;
}
\end{verbatim}

\begin{verbatim}
int max(int x,int y){
     int m=(x-y)>>31;
     return y&m|x&~m;
}
\end{verbatim}

Nếu \texttt{x} đang bằng một trong hai biến \texttt{a}, \texttt{b}, để chuyển
sang biến còn lại:
\begin{verbatim}
x^=a^b
\end{verbatim}

Đổi chỗ không dùng biến phụ:
\begin{verbatim}
void swap(int& x,int& y){x^=y;y^=x;x^=y;};
\end{verbatim}

Tính trung bình hai số nguyên không tràn:
\begin{verbatim}
int ave(int x,int y){return (x&y)+((x^y)>>1);};
\end{verbatim}

\section{Lưu ý khi dùng mảng nhiều chiều}

Bộ nhớ máy tính là tuyến tính, còn chương trình thường dùng cấu trúc nhiều
chiều như ma trận hoặc bảng quy hoạch động 3, 4 chiều. C cần ánh xạ chỉ số
nhiều chiều xuống địa chỉ một chiều. Với mảng hai chiều, các hàng được đặt
liên tiếp; với nhiều chiều hơn, có thể hiểu đệ quy: đặt \(x\) mảng con
\((d-1)\) chiều liên tiếp.

\subsection{Chi phí định địa chỉ}

Ví dụ:
\begin{verbatim}
int a[10][10][10][10][10];
int main(){
      int x=3,y=7,z=4,w=2,l=1;
      a[x][y][z][w][l]=5;
      return 0;
}
\end{verbatim}
gcc sinh mã định địa chỉ dạng:
\begin{verbatim}
movl      -20(%ebp), %ecx
movl      -24(%ebp), %ebx
movl      -28(%ebp), %esi
movl      -32(%ebp), %edx
movl      -36(%ebp), %edi
movl      %edx, %eax
sall $2, %eax
addl %edx, %eax
leal (%eax,%eax), %edx
imull     $10000, %ecx, %eax
addl %eax, %edx
imull     $1000, %ebx, %eax
addl %eax, %edx
imull     $100, %esi, %eax
leal (%edx,%eax), %eax
addl %edi, %eax
movl      $5, a(,%eax,4)
\end{verbatim}
Chỉ để làm \texttt{a[x][y][z][w][l]=5}, đã có 17 lệnh, trong đó 3 lệnh
\texttt{imul}. Vì vậy trong DP, viết:
\begin{verbatim}
if(a[x-1][yy][z-3]<a[x][y][z])a[x][y][z]=a[x-1][yy][z-3];
\end{verbatim}
có thể đắt hơn nhiều so với:
\begin{verbatim}
int &mn=a[x][y][z],v=a[x-1][yy][z-3];
if(v<mn)mn=v;
\end{verbatim}

Khi truy cập có quy luật, nên dùng con trỏ và offset đã tính sẵn. Ví dụ sau
\texttt{a[x][y][z]} thường là \texttt{a[x][y][z+1]}, địa chỉ chỉ cần tăng.
Nếu chuyển trạng thái luôn từ \texttt{a[x-1][y-2]} sang \texttt{a[x][y]},
có thể tính trước offset giữa hai vị trí.

\subsection{Hiệu quả tầng thấp của truy cập mảng}

Xét chương trình:
\begin{verbatim}
int a[2048][2048];
int main(){
      int i,x,y;
      for(i=1;i<=100;i++){
            for(x=0;x<2048;x++)
                 for(y=0;y<2048;y++)
                      a[x][y]=x+y;
      }
      return 0;
}
\end{verbatim}
Duyệt theo hàng, thời gian \texttt{user time} là 1.760s. Nếu đổi thứ tự vòng
\texttt{x} và \texttt{y}, tức duyệt theo cột, thời gian thành 18.757s, chậm
hơn hơn 10 lần. Nếu tăng kích thước mảng thành \(2049\times 2049\), dữ liệu
nhiều hơn nhưng thời gian chỉ 4.644s.

Duyệt theo hàng nhanh vì đúng thứ tự lưu trong bộ nhớ, thân thiện với cache.
Duyệt theo cột phải nhảy xa, dễ mất dữ liệu trong L1. Yếu tố này giải thích
mức 1.7s lên 4.6s. Kết quả 18s còn do kích thước 2048 là lũy thừa của 2. Khi
bước nhảy địa chỉ là lũy thừa của 2, các bit dùng để chọn nhóm cache lặp lại,
gây xung đột cache liên tục; trong trường hợp xấu, mỗi truy cập đều miss.

Kết luận thực dụng: với mảng nhiều chiều lớn, tránh để kích thước của bất kỳ
chiều nào là lũy thừa của 2; tốt nhất dùng kích thước lẻ hoặc nới thêm một
chút. DP nén trạng thái thường có một chiều \(2^k\), lúc đó tăng kích thước
mảng thêm ít phần tử đôi khi giúp đáng kể.

\section{Ví dụ ứng dụng}

Trong nhiều bài thi, thuật toán đúng là đủ lấy điểm tối đa. Nhưng khi chưa
tìm được thuật toán tối ưu, các tối ưu hằng số có thể giúp qua thêm 1 đến 4
test ở một số bài; cộng lại vẫn đáng kể.

\subsection{Số Mersenne}

Nguồn: NOIP 2003, bảng phổ cập.

Bài toán: tính \(k\) chữ số thập phân cuối của \(2^p-1\), với
\(p\le 3\,100\,000\), \(k=500\).

Thuật toán mộc mạc: lặp \(p\) lần, mỗi lần nhân số lớn với 2, độ phức tạp
\(\Oh(kp)\). Kết quả qua 7 test. Tối ưu phép nhân số lớn, loại bỏ phép chia
và nhánh điều kiện, qua đủ 10 test.

\subsection{Vũ điệu lộng lẫy}

Nguồn: NOI 2005.

Bài toán có lưới \(N\times M\), \(N,M\le 200\). Vật thể bắt đầu tại
\((X,Y)\). Có chuỗi hành động dài \(T\le 40000\), gồm các hướng đông, nam,
tây, bắc, được chia thành \(S\le 200\) đoạn, mỗi đoạn có hướng giống nhau.
Vật thể đi từng bước theo chuỗi, không được đụng chướng ngại và không ra khỏi
lưới. Hỏi cần xóa ít nhất bao nhiêu phần tử trong chuỗi để đi hợp lệ.

DP mộc mạc dùng \(F[i][j][k]\): đang ở \((j,k)\), xét từ đoạn hành động thứ
\(i\), số phần tử ít nhất cần xóa. Khi chuyển, liệt kê số phần tử \(t\) bị xóa
trong đoạn \(i\), độ phức tạp \(\Oh(NMT)\). Bản mộc mạc qua 6 test; tối ưu
định địa chỉ mảng nhiều chiều qua đủ 10 test.

\subsection{Dịch tác phẩm Maya}

Nguồn: IOI 2006.

Bài toán: cho hai xâu \(W\) độ dài \(G\le 3000\) và \(S\) độ dài
\(N\le 3000000\), bảng chữ cái \(A=52\). Đếm số lần một hoán vị của \(W\)
xuất hiện trong \(S\).

Cài đặt mộc mạc duyệt vị trí bắt đầu trong \(S\), kiểm tra tần suất ký tự của
đoạn bằng \(\Oh(A)\), tổng \(\Oh(AN)\). Bản này qua 12 test; tối ưu vòng lặp
và định địa chỉ mảng bằng hợp ngữ qua đủ 16 test.

\section{Tổng kết}

Các kỹ thuật trên rời rạc và không có ``thuốc vạn năng''. Nhưng có vài nguyên
tắc đáng giữ.

\paragraph{Tối ưu quá sớm là nguồn gốc của nhiều rắc rối.}
Câu nói nổi tiếng nhờ Donald Knuth rất thực dụng trong thi đấu. Tối ưu có
giá: mã khó đọc, khó debug, khó bảo trì, và có thể khiến ta không còn kiểm
soát được chương trình vốn có thể viết đúng. Ngoài ra, tối ưu tốn thời gian.
Khi viết bản đầu, đặc biệt trong cuộc thi, không nên thêm tối ưu làm tăng độ
khó tư duy hoặc cài đặt. Hãy có chương trình đúng trước, rồi tìm bottleneck.
Nếu một hàm chỉ chiếm 0.01 giây, tối ưu nó 100 lần cũng không ảnh hưởng đáng
kể. Thường 99\% thời gian nằm trong 1\% mã, hoặc có bài 70\% thời gian nằm ở
I/O. Hãy tối ưu đúng chỗ; với phần không phải bottleneck, cài đặt rõ ràng là
đủ.

\paragraph{Tối ưu không có điểm kết thúc.}
Sau khi tìm được bottleneck, đừng vội nói đoạn mã không thể nhanh hơn. Điều
chỉnh logic ở ngôn ngữ bậc cao, phối hợp với trình biên dịch, hoặc phối hợp
với CPU ở tầng hợp ngữ đều có không gian đào sâu. Nhưng tối ưu thật sự đòi
hỏi kiến thức và kinh nghiệm về CPU, hợp ngữ, trình biên dịch. Nó có quy luật
nhưng cũng đầy ngoại lệ. Cách duy nhất để biết tối ưu có hiệu quả là tự cài
và đo được cải thiện.

\paragraph{Keep it simple and stupid.}
Tối ưu tầng thấp thường làm mọi thứ phức tạp hơn, trái tinh thần KISS. Nhưng
khi một tối ưu trở nên quá rối, có thể ta đang đi sai hướng; nên dừng lại và
tìm cách đơn giản hơn. Nếu buộc phải viết mã xấu để đạt hiệu năng cực hạn,
hãy giới hạn nó vào đúng vài đoạn nóng. Phần còn lại của chương trình vẫn nên
giữ đơn giản, rõ ràng và dễ kiểm soát.

Bài viết chỉ chạm tới bề mặt của nhiều nguyên tắc và kỹ thuật. Hy vọng nó gợi
mở cho người đọc thấy trong lập trình còn có một vùng thú vị như vậy, và có
thể khai thác thêm trong thực chiến.

\section{Tài liệu tham khảo}

\begin{enumerate}
  \item \textit{Journey of Games: My Programming Reflections}.
  \item \textit{Professional Assembly Language}.
  \item \textit{Assembly Language for Intel-Based Computers}, Fifth Edition.
  \item \textit{Optimizing Software in C++: An Optimization Guide for
  Windows, Linux and Mac Platforms}.
  \item \textit{Michael Abrash's Graphics Programming Black Book Special
  Edition}.
  \item \textit{Coding ASM: Intel Instruction Set Codes and Cycles}.
  \item \textit{Intel 64 and IA-32 Architectures Software Developer's Manual}.
  \item \textit{Optimizing Subroutines in Assembly Language: An Optimization
  Guide for x86 Platforms}.
\end{enumerate}

\section{Lời cảm ơn}

Tác giả cảm ơn thầy Zhang Junliang vì sự hướng dẫn tận tâm; cảm ơn đàn anh
Feng Guodong (bluemary) vì nhiều năm giúp đỡ vô tư; cảm ơn giáo sư Zuo Jie
của Đại học Tứ Xuyên vì các ý kiến quý báu cho bài viết; cảm ơn mọi đồng đội
và đối thủ đã giúp tác giả trên con đường thi đấu; và cảm ơn cha mẹ cùng giáo
viên chủ nhiệm vì sự ủng hộ và bao dung.

\end{document}
